\documentclass[11pt]{article}
\usepackage[margin=1in]{geometry}
\usepackage{amsmath,amssymb,amsthm,booktabs,hyperref,enumitem}
\newtheorem{theorem}{Theorem}
\title{Counterexamples to both Bernoulli--Toda denominator claims}
\author{Principia Math In Progress}
\date{21 September 2026}
\begin{document}
\maketitle
\noindent\textbf{Task 1: MathDB 371655.}
For an even positive integer $N$, let
\[
 P(N)=\{p\text{ prime}:p-1\mid N\},\qquad
 D_N=\prod_{p\in P(N)}p,
 \qquad F(d)=\min\{N>0:N\text{ even},\ D_N=d\}.
\]
The formula for the denominator $D_N$ of the Bernoulli number $B_N$
is the classical von Staudt--Clausen theorem, used in the source
\cite[Section 5.1]{M}. Its Toda primes are
\[
 T(n)=\left\{p\text{ odd prime}:p-1\mid4n,
        \ \gcd\left(p,\frac{4n}{p-1}\right)=1\right\},
 \qquad t(n)=|T(n)|.
\]
Conjecture 5.6 of \cite{M} asserts that, for a Bernoulli denominator $d$,
(i) if $F(d)=4a$ and $D_{4am}=D_{4a}=d$, then $T(am)=T(a)$;
and (ii) if $D_{4am}\ne D_{4a}=d$, then $t(am)\geq t(a)$.
Our second example also satisfies $F(d)=4a$, whether or not that
condition is carried over to clause (ii).

\begin{theorem}
Both clauses of Conjecture 5.6 are false.
\begin{align*}
 \textup{(i)}\quad &a=23,\ m=47,\ d=1410:\qquad
 F(d)=92,\quad D_{4324}=D_{92}=1410,\\
 &T(23)=\{3,5,47\},\qquad T(1081)=\{3,5\};\\[3pt]
 \textup{(ii)}\quad &a=58,\ m=13747=59\cdot233,\ d=412410:\qquad
 F(d)=232,\\
 &D_{3189304}=192595470=467D_{232}>D_{232},\\
 &T(58)=\{3,5,59,233\},\qquad
 T(797326)=\{3,5,467\}.
\end{align*}
\end{theorem}

\paragraph{An elementary reduction.}
For an odd prime $p$ with $p-1\mid4n$, the condition
$p\mid4n/(p-1)$ is equivalent to $p\mid n$, since
$\gcd(p,p-1)=\gcd(p,4)=1$. Thus, writing $\Omega(n)$ for the prime
divisors of $n$,
\begin{equation}\label{T}
 T(n)=P(4n)\setminus\bigl(\{2\}\cup\Omega(n)\bigr).
\end{equation}
This equivalent definition already appears in the source and in
the 2018 MathOverflow discussion \cite{MO}. We use it only to simplify
the finite arithmetic.

\paragraph{Proof of (i).}
Since $92=4\cdot23$, its even divisors are $2,4,46,92$.
Adding one gives $3,5,47,93$, of which exactly the first three are prime.
An odd divisor plus one is an even number, prime only for divisor one.
Consequently $P(92)=\{2,3,5,47\}$ and $D_{92}=1410$.
If $D_N=1410$, both $5$ and $47$ divide $D_N$, so $4\mid N$ and
$46\mid N$. Hence $92=\operatorname{lcm}(4,46)\mid N$, proving
$F(1410)=92$.

Now $4324=4\cdot23\cdot47$. Relative to $92$, the additional even
divisors are $94,188,2162,4324$. Each gives a composite successor:
\[
 95=5\cdot19,\quad 189=3\cdot63,\quad
 2163=3\cdot721,\quad 4325=5\cdot865.
\]
The old composite successor is $93=3\cdot31$. These exhaust all even
divisors, so $P(4324)=P(92)$ and $D_{4324}=D_{92}$.
Equation \eqref{T}, with $1081=23\cdot47$, gives
$T(23)=\{3,5,47\}$ but $T(1081)=\{3,5\}$, disproving (i).

\paragraph{Proof of (ii).}
We have $232=8\cdot29$. Its even divisors are
$2,4,8,58,116,232$, whose successors are
$3,5,9,59,117,233$. Hence
\[
 P(232)=\{2,3,5,59,233\},\qquad
 D_{232}=2\cdot3\cdot5\cdot59\cdot233=412410.
\]
Any $N$ with this denominator has $233-1=232\mid N$.
Thus $F(412410)=232$.

Set $n=58\cdot59\cdot233=797326$, so
$4n=2^3\cdot29\cdot59\cdot233$. Every even divisor of $4n$ is
$2^j u$ with $j\in\{1,2,3\}$ and
\[
 u\in\{1,29,59,233,1711,6757,13747,398663\}.
\]
The following table lists \emph{every} such divisor's successor.
A bold entry is prime; every other entry has an explicit proper factor.
\begin{center}
\small
\begin{tabular}{rlll}
\toprule
$u$ & $2u+1$ & $4u+1$ & $8u+1$\\
\midrule
1 & $\mathbf{3}$ & $\mathbf{5}$ & $9=3\cdot3$\\
29 & $\mathbf{59}$ & $117=3\cdot39$ & $\mathbf{233}$\\
59 & $119=7\cdot17$ & $237=3\cdot79$ & $473=11\cdot43$\\
233 & $\mathbf{467}$ & $933=3\cdot311$ & $1865=5\cdot373$\\
1711 & $3423=3\cdot1141$ & $6845=5\cdot1369$ & $13689=3\cdot4563$\\
6757 & $13515=3\cdot4505$ & $27029=151\cdot179$ & $54057=3\cdot18019$\\
13747 & $27495=3\cdot9165$ & $54989=11\cdot4999$ & $109977=3\cdot36659$\\
398663 & $797327=41\cdot19447$ & $1594653=3\cdot531551$ & $3189305=5\cdot637861$\\
\bottomrule
\end{tabular}
\end{center}
For the entries declared prime, trial division by primes at most
$\sqrt{467}<22$ suffices; none of $2,3,5,7,11,13,17,19$ divides
$467$, and the same test up to the appropriate square root proves
primality of $59$ and $233$. The odd divisors contribute only prime 2.
We have therefore proved
\[
 P(3189304)=\{2,3,5,59,233,467\},\qquad
 D_{3189304}=467D_{232}=192595470.
\]
Since $\Omega(58)=\{2,29\}$ and
$\Omega(797326)=\{2,29,59,233\}$, equation \eqref{T} now gives
the two claimed Toda sets. In particular $t(797326)=3<4=t(58)$,
disproving (ii), even with the minimum-index condition imposed.
\hfill$\square$

\paragraph{Scope and verification.}
These examples disprove both clauses of the stated denominator
conjecture. They do not refute the separate conjecture $t(n)\geq2$
for every positive integer, or the source's restricted conjecture
about multiples of 3. All displayed numbers and hypotheses are exact;
no transcription or omitted-hypothesis issue is used.
The elementary characterization of Toda primes, von Staudt--Clausen,
and minimum-index divisibility are classical or source material,
not claimed as new.

The verifier checks all divisor successors and proper factors,
independently sieves every prime up to the largest possible candidate
and applies the original gcd definition, checks both minimum indices,
and computes three smaller-index Bernoulli numbers exactly with FLINT.
The large-index denominator is certified by the exhaustive prime
and divisor calculations, without computing its enormous numerator.
All 15 pages of the current target source were read; the author's
15-page copy has identical extracted text except for the arXiv stamp.
No target journal final or earlier full target claim was identified
in the recorded bounded audit. Review is by the authoring agent only.

\begin{enumerate}[label=\textbf{Task \arabic*:},start=2,leftmargin=*]
\item Classify the Bernoulli denominators $d$ with $4\mid F(d)$ for
which the Toda set is constant on all indices $N$ satisfying $D_N=d$.
\item Determine whether clause (i) has infinitely many counterexamples
with both $a$ and $m$ prime.
\item Determine whether $t(am)-t(a)$ is unbounded below when
$F(D_{4a})=4a$ and $D_{4am}>D_{4a}$.
\item Find sufficient arithmetic conditions on $a,m$ that recover
$t(am)\geq t(a)$ despite the loss of primes dividing $m$.
\end{enumerate}
\small
\begin{thebibliography}{9}
\raggedright
\bibitem{M} S. McKean, Toda primes,
\href{https://arxiv.org/abs/2511.19744}{arXiv:2511.19744v1} (2025),
Definition 1.1; Conjecture 5.6, p.~8.
\bibitem{PW} C. Pomerance and S. S. Wagstaff Jr.,
The denominators of the Bernoulli numbers,
Acta Arith. 209 (2023), 1--15,
\href{https://doi.org/10.4064/aa210601-16-11}{doi:10.4064/aa210601-16-11}.
Proposition 1 gives the general identity $F(d)=\operatorname{lcm}_{p\mid d}(p-1)$.
\bibitem{MO} A. Milovanov, Number theory question from Homotopy groups
of spheres, MathOverflow (15 April 2018),
\href{https://mathoverflow.net/q/297915}{question 297915},
including G. Martin's answer.
\end{thebibliography}
\end{document}
